% recall-drill-named-facts.tex — the exact names, numbers and lists of iOS / Swift app development:
% layout priorities, stack-view enums, URLError codes, DecodingError cases, callback and state orders,
% responder chain vs hit-testing, memory rules, SwiftUI wrappers, test doubles, GoF intents, SOLID.
% Sources (checked 2026-10-09):
%   UIKit SDK header NSLayoutConstraint.h (iOS 13 SDK, mirrored at github.com/xybp888/iOS-SDKs:
%     required 1000, defaultHigh 750, defaultLow 250, fittingSizeLevel 50; Catalyst 510 / 500 / 490;
%     priority may not change from / to required once added — exception)
%   developer.apple.com/library/archive …/AutolayoutPG: AnatomyofaConstraint (priority 1 … 1000; cluster at
%     250 / 500 / 750 / 1000, one or two points off to prevent ties; hugging 250, compression 750, easier
%     to stretch than shrink) · ConflictingLayouts (breaks constraints until valid; symbolic breakpoint
%     UIViewAlertForUnsatisfiableConstraints; failure point at 999; no required hugging / compression)
%   developer.apple.com/documentation/uikit: uistackview (axis default horizontal; distribution + alignment
%     cases, default fill; spacing 0, minimum for equalSpacing / equalCentering, negative overlaps;
%     spacingUseSystem iOS 11) · uiview/translatesautoresizingmaskintoconstraints (true in code, false from
%     IB; leave alone on arrangedSubviews, cells, a VC's view) · uitableview/rowheight + estimatedrowheight
%     (default automaticDimension; IB cells; estimate must not be 0; heightForRowAt runs for every row,
%     1000 or more rows) · uiuserinterfacesizeclass · uiviewcontroller/loadview() (never call, no super) ·
%     viewisappearing(_:) (iOS 13; after added to hierarchy and laid out by superview; traits + geometry
%     valid; layout callbacks may run many times) · uiapplication/state · uiscene/activationstate ·
%     uiscenedelegate · managing-your-app-s-life-cycle (unattached, background, suspended, disconnect) ·
%     using-responders-and-the-responder-chain-to-handle-events (next responders; app delegate only if a
%     UIResponder; gesture recognizers first; nil-target actions) · uiview/hittest(_:with:) (hidden,
%     interaction disabled, alpha < 0.01 ignored; outside bounds never hit, even unclipped)
%   developer.apple.com/videos/play/wwdc2023/10055 (viewIsAppearing back-deploys to iOS 13)
%   developer.apple.com/documentation/technotes/tn3187 (scene opt-in via UIApplicationSceneManifest; log
%     from iOS 18.4; delegate method map) · ios-ipados-release-notes/ios-ipados-27-release-notes (UIKit
%     deprecations: apps built with the latest SDK must adopt the scene-based life cycle or fail to launch)
%   developer.apple.com/documentation/foundation: jsondecoder keyDecodingStrategy / convertFromSnakeCase
%     (rules; base_uri -> baseUri; use CodingKeys for acronyms) · dateDecodingStrategy (default
%     deferredToDate; the six cases) · urlsession (tasks start suspended, resume(); delegateQueue: all
%     delegate calls and completion handlers, nil -> session-created serial queue; delegate held strongly
%     until invalidated; shared session limits; data(from:) iOS 15) · urlsessionconfiguration (default /
%     ephemeral / background; request timeout 60 s reset by data; resource 7 days; 6 connections per host,
%     HTTP/1.1 only) · urlsessiontaskdelegate didCompleteWithError (only client-side errors) ·
%     urlsession/authchallengedisposition (4 cases) · urlerror/code · notificationcenter addObserver
%     (selector form: no removal needed from iOS 9; block form must be removed)
%   Foundation SDK header NSURLError.h (same mirror: -999, -1000 … -1022, -1200, -1202 values)
%   github.com/swiftlang/swift-foundation FoundationEssentials/Date.swift (Codable = seconds since the 2001
%     reference date; 978307200 s after 1970)
%   developer.apple.com/documentation/swift: decodingerror (4 cases; Context: codingPath, debugDescription,
%     underlyingError) · encodingerror (invalidValue) · isknownuniquelyreferenced(_:) (strong refs only;
%     not thread-safe) · choosing-between-structures-and-classes · iteratorprotocol/next() ·
%     sequence/makeiterator()
%   github.com/swiftlang/swift-book TSPL.docc: AutomaticReferenceCounting (weak: optional var, set to nil,
%     no observers, shorter lifetime; unowned: same or longer lifetime, runtime error, unowned(unsafe)) ·
%     Expressions (capture list initialised when the closure is created; prints "0 10"; reference semantics)
%     · ClassesAndStructures (arrays, dictionaries, strings share storage until modified) · Properties
%     (stored type properties initialised once even across threads; lazy has no such guarantee) · Methods
%     (no mutating call on a let struct) · Protocols (weak delegates, class-only = AnyObject)
%   github.com/swiftlang/swift-evolution: SE-0244 (opaque result types, Swift 5.1) · SE-0390 (deinit on
%     noncopyable structs / enums, Swift 5.9; the Book still says "only class types") · SE-0327 (actor
%     deinitializers) · proposals/testing/ST-0008 (exit tests, Swift 6.2; iOS not
%     supported)
%   developer.apple.com/documentation/swiftui: state, binding, stateobject (iOS 14; created only once in the
%     container's lifetime), observedobject (no default value; not for @Observable), environmentobject,
%     bindable (17), migrating-from-the-observable-object-protocol-to-the-observable-macro (State,
%     Environment, Bindable; body-read tracking; ObservationIgnored) · observation/observable() (iOS 17) ·
%     scenephase · combine/anycancellable (cancel() on deinit) · foundation/undomanager registerUndo
%   developer.apple.com/documentation/xcode: gathering-information-about-memory-use (Debug Memory Graph,
%     Malloc Stack, Allocations + Mark Generation, simulator never warns) · understanding-hangs-in-your-app
%     (< 100 ms rarely noticeable; tools report > 250 ms) · sigkill (termination codes)
%   keith.github.io/xcode-man-pages/leaks.1.html (Xcode's leaks(1): unreferenced malloc buffers, root cycles)
%   developer.apple.com/design/human-interface-guidelines/accessibility (iOS controls 44x44 pt, min 28x28)
%   developer.apple.com/documentation/xctest/set-up-and-tear-down-state-in-your-tests (exact order; crash =
%     no teardown) · /testing (Swift 6, Xcode 16) · testing/expectations (#expect continues, #require throws
%     ExpectationFailedError; optional unwrap overload) · testing/trait · github.com/swiftlang/swift-testing
%     MigratingFromXCTest.md (parallel by default; XCTest sequential)
%   martinfowler.com/bliki/TestPyramid.html (Cohn, Succeeding with Agile, 2009; ice-cream cone) ·
%     /articles/practical-test-pyramid.html (Unit / Service / UI) · /articles/mocksArentStubs.html
%     (Meszaros' five doubles, quoted definitions) · /articles/injection.html (Jan 2004; name settled;
%     constructor / setter / interface injection)
%   testing.googleblog.com/2015/04/just-say-no-to-more-end-to-end-tests.html (70/20/10 "as a good first guess")
%   people.csail.mit.edu/addy/pattern (Gamma, Helm, Johnson, Vlissides, Design Patterns, 1994 — the book's
%     text: intents of Adapter, Decorator (explosion of subclasses), Facade, Proxy (remote, virtual,
%     protection, smart reference), Factory Method, Abstract Factory (Kit), Builder (Director), Singleton,
%     Strategy (clients must be aware), State (who defines transitions), Template Method (Hollywood
%     principle, [Swe85]), Mediator (many-to-many -> one-to-many), Command, Observer)
%   en.wikipedia.org/wiki/Design_Patterns (1994; the 23 by category — secondary, matches the book's chapters)
%   butunclebob.com/ArticleS.UncleBob.PrinciplesOfOod (R. C. Martin's one-line SRP / OCP / LSP / ISP / DIP)
%   khanlou.com/2015/01/the-coordinator (Coordinator; Application Controller)
% @labels: area=recall kind=reference platform=ios topic=ui,networking,memory,patterns,testing discussion=328
% @tags: uilayoutpriority, content-hugging, compression-resistance, uistackview, self-sizing-cells, urlerror, decodingerror, keydecodingstrategy, urlsession, viewisappearing, scene-delegate, responder-chain, hit-testing, capture-list, copy-on-write, stateobject, observable, test-pyramid, test-doubles, xctest, swift-testing, gof, proxy, solid
\documentclass[8pt]{extarticle}
\usepackage{printup-sheet}
\usepackage{array}

\newcommand\ct[1]{\texttt{#1}}
\newcolumntype{L}[1]{>{\raggedright\arraybackslash}p{#1}}
% compact section heading: the subject's blocks are short, a full \section costs a line each
\newcommand\hd[1]{\par\vspace{3pt}\noindent{\small\bfseries\color{sheetBlue}#1}\par\vspace{1pt}}
\setlength{\aboverulesep}{0.25ex}
\setlength{\belowrulesep}{0.25ex}
\tikzset{
  hd/.style={font=\bfseries\small, text=sheetBlue, anchor=west, inner sep=0pt},
  l/.style={font=\tiny, text=black!75, align=center, inner sep=1pt},
  tk/.style={thin, sheetGrey},
  vc/.style={draw=sheetBlue, fill=sheetBlue!7, rounded corners=1pt, font=\tiny, align=center,
             inner sep=0.8pt, minimum height=6mm, text width=17.4mm},
  st/.style={draw=sheetGreen!60!black, fill=sheetGreen!10, rounded corners=1pt, font=\ttfamily\tiny,
             inner sep=1.5pt, minimum height=4mm},
  ar/.style={->, thin, sheetGrey},
}

\begin{document}

\sheettitle{Named facts — the numbers, names and lists of iOS development}{recall · folio}

\oneliner{the exact values, case lists and call orders that UIKit, Foundation, SwiftUI and the
test frameworks fix by definition, and the names that patterns and principles were given by
their authors — each checked against Apple's documentation, the Swift Book or the originator.
The why behind each lives in the linked folios.}

\vspace{2pt}
\noindent\begin{tikzpicture}[sheet]
  % ===== A: layout priorities
  \node[hd] at (0,3.6) {UILayoutPriority (Float, 1 … 1000)};
  \draw[thick, sheetGrey] (0.1,1.55) -- (5.2,1.55);
  \foreach \v in {50,250,490,500,510,750,999,1000} \draw[tk] ({0.1+5.1*\v/1000},1.47) -- ({0.1+5.1*\v/1000},1.63);
  \node[l, anchor=south west, text=black] (p50) at (0,2.62) {\ct{.fittingSizeLevel} \textbf{50}};
  \node[l, anchor=south, text=black] (p750) at (3.925,2.62) {\ct{.defaultHigh} \textbf{750}};
  \node[l, anchor=south east, text=black] (p1000) at (5.25,3.08) {\ct{.required} \textbf{1000}};
  \node[l, anchor=south] (p500) at (2.65,2.17) {Catalyst scene \textbf{510 · 500 · 490}};
  \node[l, anchor=south, text=black] (p250) at (1.375,1.72) {\ct{.defaultLow} \textbf{250}};
  \draw[tk, densely dotted] (0.355,2.62) -- (0.355,1.63);
  \draw[tk, densely dotted] (3.925,2.62) -- (3.925,1.63);
  \draw[tk, densely dotted] (5.2,3.08) -- (5.2,1.63);
  \draw[tk, densely dotted] (2.65,2.17) -- (2.65,1.63);
  \foreach \v/\t in {0/0,500/500,1000/1000} \node[l, anchor=north, text=sheetGrey] at ({0.1+5.1*\v/1000},1.45) {\t};
  \node[l, anchor=north, text=sheetGreen!45!black] at (1.375,1.18) {content \textbf{hugging} default:\\resists growing};
  \node[l, anchor=north, text=sheetOrange!85!black] at (3.925,1.18) {\textbf{compression} default:\\resists shrinking};
  \node[l, anchor=north east, text=sheetRed] at (5.25,1.18) {\textbf{999}:\\failure\\point};
  \node[l, anchor=north west, align=left] at (0,0.42) {$<$\,1000 = optional; lower gives way first ·
    cluster at 250 / 500 /\\750 / 1000, nudge 1--2 to break ties · 250 $<$ 750: stretches first};
  \draw[sheetGrey!50] (5.45,3.75) -- (5.45,-0.1);
  % ===== B: view controller appearance order
  \node[hd] at (5.6,3.6) {UIViewController — appearance};
  \node[vc] (lv) at (6.5,2.85) {\ct{loadView}\\never call; no \ct{super}};
  \node[vc] (dl) at (8.35,2.85) {\ct{viewDidLoad}\\once per loaded view};
  \node[vc] (wa) at (10.2,2.85) {\ct{viewWillAppear}\\traits, geometry stale};
  \node[vc, draw=sheetGreen!60!black, fill=sheetGreen!10] (ia) at (10.2,1.85) {\ct{viewIsAppearing}\\in hierarchy, sized};
  \node[vc, draw=sheetOrange, fill=sheetOrange!10] (ly) at (8.35,1.85) {\ct{viewWill/DidLayout}\\\ct{Subviews} 0 … n×};
  \node[vc] (da) at (6.5,1.85) {\ct{viewDidAppear}};
  \node[vc] (wd) at (6.5,0.85) {\ct{viewWillDisappear}};
  \node[vc] (dd) at (8.35,0.85) {\ct{viewDidDisappear}};
  \draw[ar] (lv) -- (dl); \draw[ar] (dl) -- (wa); \draw[ar] (wa) -- (ia);
  \draw[ar] (ia) -- (ly); \draw[ar] (ly) -- (da); \draw[ar] (da) -- (wd); \draw[ar] (wd) -- (dd);
  \node[l, anchor=west, align=left] at (9.3,0.85) {shown again:\\from \ct{viewWillAppear}};
  \node[l, anchor=north west, align=left, text=sheetGreen!45!black] at (5.6,0.42)
    {\ct{viewIsAppearing}: traits + geometry valid; once per\\appearance; WWDC23, back-deploys to iOS 13};
  \draw[sheetGrey!50] (11.2,3.75) -- (11.2,-0.1);
  % ===== C: scene activation states
  \node[hd] at (11.3,3.6) {UIScene.ActivationState};
  \node[st] (un) at (12.2,2.95) {unattached};
  \node[st] (fi) at (12.2,1.65) {foregroundInactive};
  \node[st] (fa) at (15.65,1.65) {foregroundActive};
  \node[st] (bg) at (12.2,0.55) {background};
  \node[st, draw=sheetGrey, fill=white, dashed, font=\tiny] (su) at (15.45,0.55) {suspended (no case)};
  \draw[ar] (un) -- (fi);
  \node[l, anchor=west, align=left] at (12.3,2.32) {\ct{scene(\_:willConnectTo:options:)}\\a system-requested scene usually goes to background};
  \draw[ar] ([yshift=1.2pt]fi.east) -- ([yshift=1.2pt]fa.west);
  \draw[ar] ([yshift=-1.2pt]fa.west) -- ([yshift=-1.2pt]fi.east);
  \node[l, anchor=south, font=\ttfamily\tiny] at (13.95,1.9) {sceneDidBecomeActive →};
  \node[l, anchor=north, font=\ttfamily\tiny] at (13.95,1.4) {← sceneWillResignActive};
  \draw[ar] ([xshift=-6pt]fi.south) -- ([xshift=-6pt]bg.north);
  \draw[ar] ([xshift=6pt]bg.north) -- ([xshift=6pt]fi.south);
  \node[l, anchor=west, align=left, font=\ttfamily\tiny] at (12.5,0.97) {$\downarrow$ sceneDidEnterBackground\\$\uparrow$ sceneWillEnterForeground};
  \draw[ar] (bg) -- node[l, below] {no callback} (su);
  \draw[ar] (su.east) -- (16.55,0.55) |- (un.east);
  \node[l, anchor=south] at (14.75,2.98) {\ct{sceneDidDisconnect(\_:)}\\from background or suspended};
  \node[l, anchor=north west, align=left, text=sheetRed] at (11.3,0.2)
    {built with the iOS 27 SDK: no scene life cycle = no launch};
\end{tikzpicture}

\begin{multicols}{2}
\raggedright
\footnotesize\setstretch{1.0}
\setlength{\parskip}{1.5pt plus 0.5pt}

\hd{Auto Layout}
{\scriptsize\setlength{\tabcolsep}{2pt}
\begin{tabular}{@{}L{17mm}L{61mm}@{}}
\toprule
\textbf{name} & \textbf{value / rule}\\
\midrule
installed priority & changing it to or from \ct{.required} after the constraint is added throws\\
autoresizing mask & \ct{translatesAutoresizingMaskIntoConstraints}: \ct{true} for a view made in code, \ct{false} from Interface Builder; leave it alone on \ct{arrangedSubviews}, cells, a VC's \ct{view}\\
conflict & engine breaks constraints until the layout is valid, logs; breakpoint \ct{UIViewAlertForUnsatisfiable}\allowbreak\ct{Constraints}; never required hugging / compression\\
self-sizing rows & \ct{rowHeight} \emph{and} \ct{estimatedRowHeight} default to \ct{.automaticDimension}; an IB cell needs \ct{rowHeight = .automaticDimension} in code; estimate $\ne$ 0; \ct{heightForRowAt} overrides and runs for every row (Apple: hurts at 1000+)\\
size class & \ct{UIUserInterfaceSizeClass}: \ct{.compact} · \ct{.regular} · \ct{.unspecified}, per axis; set by device, window, multitasking\\
\bottomrule
\end{tabular}}

\medskip
{\scriptsize\setlength{\tabcolsep}{2pt}
\begin{tabular}{@{}L{17mm}L{12mm}L{48mm}@{}}
\toprule
\textbf{UIStackView} & \textbf{default} & \textbf{cases}\\
\midrule
\ct{axis} & \ct{.horizontal} & \ct{.horizontal} · \ct{.vertical}\\
\ct{distribution} (along axis) & \ct{.fill} & \ct{.fill} · \ct{.fillEqually} · \ct{.fillProportionally} · \ct{.equalSpacing} · \ct{.equalCentering}\\
\ct{alignment} (across) & \ct{.fill} & \ct{.fill} · \ct{.leading} (= \ct{.top}) · \ct{.firstBaseline} · \ct{.center} · \ct{.trailing} (= \ct{.bottom}) · \ct{.lastBaseline}\\
\ct{spacing} & \ct{0} & a minimum for the two \ct{equal…} kinds; negative overlaps; \ct{spacingUseSystem} (iOS 11)\\
\bottomrule
\end{tabular}}

\hd{Networking \& Codable}
\begin{tikzpicture}[sheet]
  \node[box, font=\scriptsize, text width=22mm, draw=sheetRed, fill=sheetRed!6] (t) at (0,0) {\textbf{transport}\\\ct{URLError} — thrown /\\\ct{error} non-nil};
  \node[box, font=\scriptsize, text width=22mm, draw=sheetOrange, fill=sheetOrange!8] (h) at (2.7,0) {\textbf{HTTP status}\\\textbf{not} an error: check\\\ct{statusCode} in \ct{200..<300}};
  \node[box, font=\scriptsize, text width=22mm, draw=sheetBlue] (d) at (5.4,0) {\textbf{body}\\\ct{DecodingError}\\4 cases};
  \draw[flow] (t) -- (h); \draw[flow] (h) -- (d);
\end{tikzpicture}\par
{\scriptsize URLSession reports only client-side errors; a 404 arrives as a response
(\ct{as? HTTPURLResponse}).}

\smallskip
{\scriptsize\setlength{\tabcolsep}{1.6pt}
\begin{tabular}{@{}r@{\ }L{30mm}r@{\ }L{34mm}@{}}
\toprule
\multicolumn{4}{@{}l}{\textbf{\ct{URLError.Code}} (\ct{NSURLErrorDomain})}\\
\midrule
\ct{-999} & \ct{cancelled} & \ct{-1009} & \ct{notConnectedToInternet}\\
\ct{-1000} & \ct{badURL} & \ct{-1011} & \ct{badServerResponse}\\
\ct{-1001} & \ct{timedOut} & \ct{-1020} & \ct{dataNotAllowed} (cellular refused)\\
\ct{-1003} & \ct{cannotFindHost} & \ct{-1022} & \ct{appTransportSecurityRequires}\allowbreak\ct{SecureConnection}\\
\ct{-1004} & \ct{cannotConnectToHost} & \ct{-1200} & \ct{secureConnectionFailed}\\
\ct{-1005} & \ct{networkConnectionLost} & \ct{-1202} & \ct{serverCertificateUntrusted}\\
\bottomrule
\end{tabular}}

\smallskip
{\scriptsize
\textbf{\ct{DecodingError}}: \ct{.typeMismatch(type, ctx)} · \ct{.valueNotFound(type, ctx)} (null for
a non-optional) · \ct{.keyNotFound(key, ctx)} · \ct{.dataCorrupted(ctx)}; \ct{Context} =
\ct{codingPath}, \ct{debugDescription}, \ct{underlyingError}. \ct{EncodingError}: only
\ct{.invalidValue}.\par
\textbf{Keys}: \ct{CodingKeys} enum (per key) or \ct{keyDecodingStrategy}: \ct{.useDefaultKeys}
(default) · \ct{.convertFromSnakeCase} — capitalise after each \ct{\_}, drop inner \ct{\_}, keep
leading / trailing ones: \ct{base\_uri} → \ct{baseUri}, never \ct{baseURI} · \ct{.custom}.\par
\textbf{Dates}: default \ct{.deferredToDate} = \ct{Date}'s own Codable: a \ct{Double}, seconds since
\textbf{1 Jan 2001} (978\,307\,200\,s after 1970) — Unix time needs \ct{.secondsSince1970};
also \ct{.millisecondsSince1970} · \ct{.iso8601} · \ct{.formatted} · \ct{.custom}.\par}

\smallskip
{\scriptsize\setlength{\tabcolsep}{2pt}
\begin{tabular}{@{}L{15mm}L{63mm}@{}}
\toprule
\textbf{URLSession} & \\
\midrule
start & every task is created suspended → \ct{resume()}\\
callback queue & delegate calls + completion handlers on \ct{delegateQueue}; \ct{nil} → a serial queue the session creates, not main\\
async & \ct{data(from:)} (iOS 15) → \ct{(Data, URLResponse)}\\
configs & \ct{.default} (disk cache, keychain) · \ct{.ephemeral} (RAM only) · \ct{.background(withIdentifier:)} (system process; force-quit cancels)\\
numbers & request timeout \textbf{60\,s}, reset by each arriving byte · resource \textbf{7 days} · \textbf{6} connections per host, HTTP/1.1 only\\
ownership & delegate held \textbf{strongly} until \ct{invalidate…} — else a leak; \ct{.shared}: no delegate, no custom config, no background\\
challenge & \ct{urlSession(\_:didReceive:}\allowbreak\ct{completionHandler:)} → \ct{.useCredential} · \ct{.performDefaultHandling} · \ct{.cancelAuthentication}\allowbreak\ct{Challenge} · \ct{.rejectProtectionSpace}\\
\bottomrule
\end{tabular}}

\hd{Lifecycle enums \& kill codes}
{\scriptsize
\ct{UIApplication.State} and SwiftUI \ct{ScenePhase}: \ct{.active} · \ct{.inactive} ·
\ct{.background}. Scenes opt in via \ct{UIApplicationSceneManifest} (iOS 13); UIKit logs a message
without it from iOS 18.4. TN3187: \ct{applicationDidBecomeActive} → \ct{sceneDidBecomeActive}, likewise
\ct{WillResignActive}, \ct{DidEnterBackground}, \ct{WillEnterForeground}. \textbf{Hang}: most Apple
tools report a main run loop busy $>$ \textbf{250\,ms}; $<$ 100\,ms is rarely noticed. \ct{0x8badf00d} watchdog
· \ct{0xdead10cc} file / SQLite lock held while suspending · \ct{0xc00010ff} thermal ·
\ct{0xbaadca11} VoIP push without a CallKit report.\par}

\hd{SOLID — R. C. Martin's own one-liners}
{\scriptsize
\textbf{S} A class should have one, and only one, reason to change. · \textbf{O} You should be able
to extend a classes behavior, without modifying it. · \textbf{L} Derived classes must be
substitutable for their base classes. · \textbf{I} Make fine grained interfaces that are client
specific. · \textbf{D} Depend on abstractions, not on concretions.\par}

\columnbreak

\hd{Touches: down to find, up to handle}
\begin{tikzpicture}[sheet]
  \foreach \n/\x/\t in {a/0/view, b/1.0/superview, c/2.05/VC root, d/3.0/VC, e/3.9/window, f/4.95/UIApplication, g/6.2/app delegate}
    \node[st, font=\tiny, draw=sheetBlue, fill=sheetBlue!7] (\n) at (\x,0) {\t};
  \foreach \p/\q in {a/b,b/c,c/d,d/e,e/f,f/g} \draw[ar] (\p) -- (\q);
  \node[l, anchor=south west] at (-0.4,0.2) {\textbf{responder chain →} (unhandled event moves up)};
  \node[l, anchor=north west, align=left] at (-0.4,-0.2) {\textbf{hit-testing ←} from the window down, frontmost subview first,
    to the deepest view holding the point};
\end{tikzpicture}\par
{\scriptsize A presented VC's next responder is the presenting VC; the app delegate joins only if it
is a \ct{UIResponder}. \ct{hitTest(\_:with:)} skips \ct{isHidden}, \ct{isUserInteractionEnabled = false},
\ct{alpha < 0.01}; a point outside a view's bounds never hits its subviews, even with
\ct{clipsToBounds = false}; transparent pixels still hit. Gesture recognizers get touches before the
view; a \ct{nil}-target action searches the chain.\par}

\hd{Memory \& value semantics}
{\scriptsize\setlength{\tabcolsep}{2pt}
\begin{tabular}{@{}L{12mm}L{32mm}L{33mm}@{}}
\toprule
 & \ct{weak} & \ct{unowned}\\
\midrule
form & always \ct{var}, Optional & non-optional or Optional\\
target dies & ARC sets \ct{nil} (no observers fire) & access = runtime error; \ct{unowned(unsafe)} is unchecked\\
use when & the other lives \textbf{shorter} & the other lives as long or \textbf{longer}\\
\bottomrule
\end{tabular}}
{\scriptsize
\ct{var a = 0, b = 0; let c = \{ [a] in print(a, b) \}; a = 10; b = 10; c()} prints \textbf{0 10}: list
entries are copied when the closure is \emph{created}; a class reference copies the reference.
\ct{weak var delegate: P?} needs \ct{protocol P: AnyObject}.\par
\textbf{Copy-on-write}: arrays, dictionaries, strings share storage until a copy is modified; own
structs only by hand — \ct{isKnownUniquelyReferenced(\_:)} counts strong refs only, needs
synchronisation. \ct{static let}: lazy, initialised once even across threads; \ct{lazy var}: no
such guarantee. No \ct{mutating} call on a \ct{let} struct. \ct{deinit}: classes, actors and
\ct{\textasciitilde Copyable} structs / enums (SE-0390, Swift 5.9). Apple: structs by default;
classes for Objective-C interop or controlled identity.\par
\textbf{Tools}: Debug Memory Graph (+ scheme Malloc Stack for backtraces) · Allocations, Mark
Generation around one feature · \ct{leaks}: unreferenced malloc buffers + root cycles · the simulator
never warns or kills for memory.\par}

\hd{SwiftUI state}
{\scriptsize\setlength{\tabcolsep}{2pt}
\begin{tabular}{@{}L{17mm}L{5mm}L{55mm}@{}}
\toprule
\textbf{wrapper} & \textbf{iOS} & \textbf{rule}\\
\midrule
\ct{@State} & 13 & owned; also holds an \ct{@Observable} model\\
\ct{@Binding} & 13 & read / write a source of truth elsewhere (\ct{\$x})\\
\ct{@StateObject} & 14 & created \textbf{once} per container lifetime\\
\ct{@ObservedObject} & 13 & an input — never a default value; not for \ct{@Observable}\\
\ct{@Observable} & 17 & macro; a view updates only for properties \ct{body} reads\\
\ct{@Bindable} & 17 & bindings into an \ct{@Observable}\\
\bottomrule
\end{tabular}}
{\scriptsize Migration: \ct{StateObject} → \ct{State} · \ct{EnvironmentObject} → \ct{Environment} +
\ct{.environment(\_:)} · \ct{ObservedObject} → plain property · \ct{@Published} → nothing
(\ct{@ObservationIgnored} opts out). \ct{some View}: opaque result type, SE-0244, Swift 5.1.\par}

\hd{Testing}
{\scriptsize \textbf{Pyramid} (Cohn, \emph{Succeeding with Agile}, 2009): Unit · Service · UI; inverted =
ice-cream cone. \textbf{70/20/10} is Google Testing Blog 2015, ``as a good first guess'' — not Cohn.\par}
{\scriptsize\textbf{Test doubles} (Meszaros, \emph{xUnit Test Patterns}, 2007): \textbf{Dummy} passed
around, never used — fills parameter lists · \textbf{Fake} working shortcut implementation (an
in-memory database) · \textbf{Stub} canned answers to calls made during the test · \textbf{Spy} a stub
that also records how it was called · \textbf{Mock} pre-programmed expectations = the specification of
the calls.\par}
{\scriptsize\textbf{XCTest}: class \ct{setUp()} → per test: \ct{setUp() async throws} → \ct{setUpWithError()}
→ \ct{setUp()} → test → teardown blocks (LIFO) → \ct{tearDown()} → \ct{tearDownWithError()} →
\ct{tearDown() async throws} → class \ct{tearDown()}; a crash skips teardown. \textbf{Swift Testing}
(Swift 6, Xcode 16): \ct{\#expect} records and continues; \ct{\#require} throws
\ct{ExpectationFailedError}, \ct{try \#require(opt)} unwraps; parallel by default (XCTest: one by
one); traits \ct{.disabled} \ct{.enabled(if:)} \ct{.bug} \ct{.tags} \ct{.timeLimit} \ct{.serialized};
exit tests Swift 6.2, not on iOS.\par}

\hd{Patterns \& principles — the authors' words}
{\scriptsize \textbf{GoF} (Gamma, Helm, Johnson, Vlissides, 1994) — 23: \textbf{5 creational} Abstract
Factory, Builder, Factory Method, Prototype, Singleton · \textbf{7 structural} Adapter, Bridge,
Composite, Decorator, Facade, Flyweight, Proxy · \textbf{11 behavioural} Chain of Responsibility,
Command, Interpreter, Iterator, Mediator, Memento, Observer, State, Strategy, Template Method,
Visitor.\par}
{\scriptsize\setlength{\tabcolsep}{2pt}
\begin{tabular}{@{}L{12mm}L{66mm}@{}}
\toprule
Adapter & \emph{converts} an interface into the one clients expect\\
Decorator & same interface, \emph{adds} responsibilities at run time — vs ``an explosion of subclasses to support every combination''\\
Facade & \emph{one higher-level} interface over a subsystem\\
Proxy & surrogate that \emph{controls access}: remote · virtual (create on demand) · protection · smart reference\\
\midrule
Strategy & interchangeable algorithms; the \textbf{client} picks, so it must know them\\
State & object ``will appear to change its class''; transitions in the Context or (more flexible) in the states\\
\midrule
Factory M. & subclasses decide which class to instantiate — one product\\
Abstract F. & families of related objects, no concrete classes named (``Kit'')\\
\bottomrule
\end{tabular}}
{\scriptsize Template Method = ``Hollywood principle: Don't call us, we'll call you'' · Mediator:
many-to-many → one-to-many · Command → queue, log, \textbf{undo} (\ct{UndoManager.registerUndo}) ·
Observer: \ct{NotificationCenter} selector observers need no removal since iOS 9, block observers do;
\ct{AnyCancellable} cancels on deinit · Builder's \textbf{Director} · Iterator =
\ct{IteratorProtocol.next()} (\ct{nil} = end) + \ct{Sequence.makeIterator()} · Coordinator
(Khanlou, 2015) · ``Dependency Injection'': name settled in Fowler's Jan 2004 article (constructor /
setter / interface injection) — a technique; DIP is a principle.\par}

\end{multicols}

\noindent{\footnotesize\color{sheetGrey}\textit{Related:} Auto Layout · URLSession \& networking ·
Codable in depth · App \& view-controller lifecycle · UIKit touches · ARC · Value vs reference ·
SwiftUI state \& data flow · Testing fundamentals · Test doubles · XCTest \& Swift Testing ·
GoF: creational + structural · GoF: behavioural patterns · SOLID · Crashes \& symbolication}

\end{document}
